\chapter{Prologue}\label{chap:introduction}
The \emph{quantum many-body problem} stands as one of the central problems of physics of the last century. It concerns the description of the \emph{collective behaviour} of a \emph{macroscopic} number of quantum constituents. Despite the relative simplicity of the one-particle Schr\"odinger equation, the \emph{exponential growth} of the many-body Hilbert space poses a fundamental challenge to analytical and computational approaches, rendering exact solutions essentially intractable for all but the smallest system sizes.

The significance of the many-body problem can hardly be overstated as it lies at the heart of quantum field theory, condensed matter physics, quantum chemistry and quantum computation. Any progress on this front therefore carries implications across a wide spectrum of fields.

A crucial observation, however, is that the typical states of physical interest in the many-body Hilbert space occupy only a \emph{low-dimensional manifold}. A physically motivated approach to the many-body problem should therefore aim to efficiently parametrise this tiny low-complexity manifold and develop methods to project and diagonalise the Hamiltonian within it. This philosophy exactly underlies the \emph{tensor network} approach, which forms the central technical toolbox of this thesis. In essence, tensor network states provide efficient parametrisations of states satisfying an \emph{area law} for their \emph{entanglement entropy}.
\section{The exponential wall}
The hallmark of a quantum many-body system is the \emph{tensor product structure} of its Hilbert space. By this, it is meant that the full many-body Hilbert space $\mc H$ can be written as a tensor product of constituent spaces hosting localised \emph{quantum modes} or \emph{particles}:
\begin{equation}\label{eq:HilbertSpace}
	\mc H = \bigotimes_{\msf i=1}^N \mc H_\msf i.
\end{equation}
Throughout, we will assume the local Hilbert space dimension $d_\msf i:=\dim_\mbb C\mc H_\msf i$ to be finite, which is naturally true in many cases of interest, for example when $\mc H_\msf i$ is the Hilbert space of a qubit $\mc H_\msf i\simeq\mbb C^2$. Even when $\mc H_\msf i$ is inherently infinite-dimensional such as in the case of bosonic modes, a truncation of the local basis is often permitted for many physically relevant systems.

Provided an (orthonormal) basis for the local modes, $\mc H_\msf i:= {\rm Span}_{\mbb C}\{ | i\ra _{\msf i} \ |\  i=1,2,\dots, d_\msf i \}$, let us consider the (orthonormal) tensor product basis for $\mc H$ denoted by
\begin{equation}
	| i_1, i_2, \dots , i_N \ra := | i_1 \ra \otimes | i_2 \ra \otimes \dots \otimes | i_N \ra.
\end{equation}
In this basis, a many-body wave function $|\Psi\ra$ can be expressed as
\begin{equation}
	| \Psi \ra = \sum_{ \{i_j\} } \Psi_{i_1,i_2,\dots,i_N} | i_1, i_2, \dots , i_N \ra,
\end{equation}
with amplitudes $\Psi_{i_1,i_2,\dots,i_N}:=\la i_1, i_2, \dots , i_N | \Psi \ra $. Notably, a generic wave function is specified by a number of complex coefficients $\Psi_{i_1,i_2,\dots,i_N}$ that scales \emph{exponentially} in the number of constituent particles. Consequently, an exact treatment of even a moderate number of particles very quickly becomes computationally intractable. However, as argued below, most many-body states cannot be created by any reasonable physical mechanism due to the \emph{superposition principle} of quantum mechanics, which leads to the existence of highly \emph{entangled} many-body states. 

\bigskip\noindent A precise mathematical characterisation of entanglement is as follows. Consider a \emph{bipartition} of the Hilbert space in a subsystem $A$ and its complement $\bar A$, i.e. $\mc H=\mc H_A\otimes \mc H_{\bar{A}}$, and write $d_A=\dim(\mc H_A)$,  $d_{\bar{A}}=\dim(\mc H_{\bar{A}})$. By the \emph{singular value decomposition} (SVD), any state in $\mc H$ can be written in the form
\begin{equation}\label{eq:SchmidtState}
	\ket{\Psi} = \sum_{k=1}^r s_k | \psi_k^A \ra \otimes | \psi_k^{\bar{A}} \ra,
\end{equation}
where $| \psi_k^A \ra$ and $| \psi_k^{\bar{A}} \ra$, $r\leq\min(d_A,d_{\bar{A}})$, are orthonormal bases for $\mc H_A$ and $\mc H_{\bar{A}}$ respectively. The non-negative numbers $s_k$ for $k=1,2,...,r$ are called \emph{singular values} or \emph{Schmidt coefficients}. The state $| \Psi \ra$ is said to be \emph{entangled} between the parties $A$ and $\bar A$ when the Schmidt rank $r$ is strictly larger than $1$, and to be \emph{maximally entangled} when all singular values are equal (and non-zero). The state $| \Psi \ra$ is thus \emph{unentangled} if and only if it is a (tensor) product state, viz. $| \Psi \ra = | \psi^A \ra \otimes | \psi^{\bar A} \ra$.

Information about the entanglement can be compressed into one number, the \emph{entanglement entropy} $S_A$, which is defined as the \emph{von Neumann entropy} of the reduced density matrix $\rho_A$ (or equivalently $\rho_{\bar A}$), and which can simply be expressed as:
\begin{equation}\label{eq:ReducedRho}
	\rho_A = {\rm Tr}_{\bar A}(| \Psi \ra\la \Psi |) = \sum_k s_k^2 | \psi_k^A \ra\la \psi_k^A |,
\end{equation}
such that:
\begin{equation}
	S_A = -{\rm Tr}(\rho_A\log(\rho_A)) = -\sum_k s_k^2 \log(s_k^2).
\end{equation}
%Two of the central results in entanglement theory state that any pair of bipartite states which have the same amount of entanglement, as measured by the entanglement entropy, can be converted into each other by means of only local operations on either subsystem, supplemented with classical communication; moreover any (partially) entangled bipartite state can be created from a maximally entangled resource state such as a spin singlet~\cite{Bennett1995}.

More refined information about the entanglement structure of the state is encoded into the \emph{entanglement spectrum}, defined as the eigenvalues $\{s_k^2\}_k$ of the reduced density matrix $\rho_A$. A central result in the field of quantum information theory due to Nielsen stipulates that by making use of \emph{local operations supplemented with classical communication}, so-called LOCC transformations, a bipartite state $|\Psi_1\ra$ can be converted \emph{deterministically} into a state $|\Psi_2\ra$ if and only if the entanglement spectrum of $|\Psi_1\ra$, $\{s_{k,1}^2\}_k$, \emph{majorises} that of $|\Psi_2\ra$, $\{s_{k,2}^2\}_k$, meaning that $\sum_{k=1}^m s_{k,1}^2 \geq \sum_{k=1}^m s_{k,2}^2$, for all $m=1,\dots, r$, and where the Schmidt coefficients are assumed to be \emph{decreasingly} ordered~\cite{Nielsen1999}.

For many cases of interest --- notably the low-energy states of \emph{local gapped} Hamiltonians --- the rapid decay of the entanglement spectrum is precisely what permits an efficient \emph{tensor network} parametrisation of the state $| \Psi \ra$. As they form the central technical tool of this thesis, \cref{ch:TN} is devoted to introduce and discuss them.

From the expression \cref{eq:ReducedRho} we can also conclude that $\rho_A$ only describes a pure state, i.e. $\rank(\rho_A)=1$, when $| \Psi \ra$ is a product state; or, in other words, that the subsystem $A$ behaves as a classical mixture if the state $| \Psi \ra$ is entangled. This neatly illustrates Schr\"odinger's observation that \emph{[...] the best possible knowledge of a whole does not necessarily include the best possible knowledge of all its part [...]}~\cite{Schrodinger2008}.

\bigskip\noindent Having defined the notion of entanglement, we can make precise the above claim that most states in the quantum many-body Hilbert space are physically inaccessible. Focusing on the Hilbert space of $N$ qubits, let us consider to this end the following question put forward in ref.~\cite{Poulin2011}: what portion of the many-body Hilbert space can be prepared and explored by means of a physical mechanism, specifically by time-evolution under a local time-dependent Hamiltonian for a time that scales at most \emph{polynomially} in the system size? Following ref.~\cite{Poulin2011}, any such Hamiltonian is approximated to arbitrary precision by a polynomial-depth \emph{quantum circuit} consisting of few-qubit gates drawn from a \emph{universal gate set}, and \emph{$\varepsilon$-balls} around the outputs of these circuits are considered. In this context, an $\varepsilon$-ball around a certain state $| \Psi_0 \ra$ is defined as the collection of all states $|\Psi \ra$ such that
\begin{equation}
	|\la \Psi | \Psi_0\ra| \geq 1 - \varepsilon.
\end{equation}
A straightforward counting argument demonstrates that such balls occupy a fraction of the many-body Hilbert space that scales as $\mc O(\varepsilon^{2^N})$, i.e. as a double exponential!

From these considerations we must, perhaps quite shockingly, conclude that the vast majority of states in the Hilbert space should be deemed \emph{unphysical} in any operationally meaningful sense. Indeed, there is an overwhelming time barrier to prepare a generic state consisting of any more than a handful of constituent particles.

\bigskip\noindent Any physically motivated approach to the many-body problem should therefore consist of identifying the relevant low-complexity \emph{manifold of states} and developing the techniques needed to deal with them. Perhaps the most well-known such manifold is that of \emph{Slater determinants} of fermions. This parametrisation and its associated self-consistent Hartree-Fock method has been extremely successful for almost-free fermion systems and has elucidated the structure of the periodic table of elements~\cite{Slater1929,Slater1935,Hartree1947}. However, when correlations are strong, perturbative approaches such as Hartree-Fock fail to capture the relevant physics. It turns out that for systems satisfying an \emph{area law} of entanglement entropy, the natural low-complexity manifold is precisely that of tensor network states. In particular, such an area law can be proven to hold for one-dimensional \emph{local gapped} \emph{lattice models}, and is observed in higher dimensions.
\section{Quantum lattice models}
\emph{Quantum lattice models} appear ubiquitously as an effective, compressed description of the low-energy physics of a condensed matter system or as spatial discretisation of a quantum field theory. In the context of this thesis, we will roughly think of a \emph{lattice} $\Lambda$ as a discretisation of the spatial manifold on which the many-body Hilbert space is supported. It consists of a collection of vertices, edges, plaquettes, \dots that we will denote by respectively $\msf V(\Lambda)$, $\msf E(\Lambda)$, $\msf P(\Lambda)$, \dots

Provided a lattice, the many-body Hilbert space \cref{eq:HilbertSpace} is constructed by assigning a local Hilbert space to the vertices, edges, plaquettes or combinations thereof. The model itself is then defined through a microscopic Hamiltonian $H$ consisting of \emph{geometrically local} terms so that it can be written as $H=\sum_\msf i h_{\msf i,K}$ where $h_{\msf i, K}$ denotes a term acting non-trivially on at most $K$ sites in the neighbourhood of site $\msf i$. Moreover, such Hamiltonians typically depend on a small number of tunable parameters that are meant to mimic among other things external magnetic fields and chemical potentials.

Throughout this thesis we are primarily interested in the \emph{ground state subspace} and \emph{low-lying excited states} of local lattice models in \emph{one} and \emph{two spatial dimensions}. Our focus on the low-energy spectrum is physically motivated by the observation that the Fermi temperature of valence electrons in materials of interest is of the order $10^4K$, well above the temperature scale of experiments. Also in the perturbative computation of cross sections in quantum field theory, the ground state is usually the reference state on top of which excitations are created.

\bigskip\noindent Finding the (approximate) ground state of a generic quantum lattice model is notoriously difficult due to the impossibility to minimise the energy of all Hamiltonian terms simultaneously~\cite{Verstraete2006a}. This phenomenon is referred to as \emph{quantum frustration} and is intimately related to the concept of \emph{monogamy of entanglement}. As a paradigmatic example, consider the antiferromagnetic spin 1/2 Heisenberg chain:
\begin{equation}
	H = \sum_{\msf i} \vec S_{\msf i} \cdot \vec S_{\msf i+1},
\end{equation}
with $\vec S=(S^x,S^y,S^z)$ the $\mathfrak{su}(2)$ generators in the spin 1/2 representation. On two sites, the unique ground state is the (maximally entangled) spin singlet $\frac{1}{\sqrt 2}( | 01 \ra - | 10 \ra)$. One might therefore hope that the ground state on any chain length is one for which all reduced density matrices $\rho_{\msf i,\msf i+1}$ are equal to a spin singlet. Such a state would indeed clearly minimise all exchange terms $\vec S_\msf i\cdot\vec S_{\msf i+1}$. However, such a state does not exist in virtue of \emph{monogamy of entanglement} which stipulates that a maximally entangled pair of particles cannot be entangled with a third one~\cite{Coffman1999,Osborne2006}. Instead, each spin must optimally distribute its entanglement among its neighbours, and it is precisely this frustration that gives rise to a strongly correlated ground state characterised by intricate \emph{entanglement patterns}.

\bigskip\noindent The nature of these entanglement structures depends strongly on whether the model is gapped. A model is said to posses an \emph{(energy) gap} if the excitation energy of the lowest-lying excited state(s) is non-vanishing in the thermodynamic limit $N\rightarrow\infty$.\footnote{A state which has a finite energy gap above the true ground state for any finite system size, but whose energy splitting vanishes exponentially in the thermodynamic limit $N\mapsto\infty$ is thus also a ground state.} Such a system is characterised by exponentially decaying \emph{connected correlation functions} to which a finite \emph{correlation length} $\xi$ can be associated, viz.:
\begin{equation}
	\la O(\vec r) O(\vec 0)\ra_{\rm conn.} \sim \exp(-|\vec r|/\xi),
\end{equation}
for a local operator $O$~\cite{Hastings2004,Hastings2005}. The entanglement structure of gapped systems differs vastly from that of arbitrary states in the many-body Hilbert space. By virtue of the \emph{locality} of physical interactions, gapped ground states obey an \emph{area law} of their entanglement entropy. In generic spatial dimensions, the area law of entanglement entropy stipulates that the von Neumann entropy between any region $A$ and its complement does not scale as the volume of its interior (as could be expected for generic quantum states)  but rather by the surface area of its boundary $\partial A$. Specialising to the case of one-dimensional chains, the area of a contiguous region is simply a constant and for this case the area law was formally proven by Hastings by a clever application of Lieb-Robinson bounds~\cite{Hastings2007,Lieb1972}. In fact, it can be shown that the entanglement entropy scales with the correlation length according to $\log(\xi)$~\cite{Calabrese2004}. \footnote{The entanglement entropy of \emph{critical} systems mildly violates the area law, since in that case the entanglement entropy of a contiguous region of $L$ spins scales as $\log(L)$.} In higher dimensions, the area law is not formally proven, but observed to hold in important cases of interest.

The low-lying spectrum on top of gapped ground states likewise exhibit a characteristic structure.
Elementary excitations are typically organised in a few well-defined isolated branches of \emph{quasiparticle} excitations with well-defined energy and momentum~\cite{Bijl1941,Feynman1953,Feynman1954,Haegeman2011b,Vanderstraeten2013}. At higher energies, quasiparticle combine into \emph{scattering states} whose energy is bounded from below by the two-particle excitation energy.
\section{Symmetries and quantum phases of matter}\label{eq:SymGappedPhases}
The central challenge in the many-body problem is to determine and characterise the \emph{phase diagram} of a given microscopic model, provided a \emph{symmetry} of the model. 

A symmetry is, roughly speaking, a collection of transformations that leaves invariant the dynamics of the model. A standard distinction is made between \emph{internal} and \emph{spacetime} symmetries. Spacetime symmetries act on the underlying lattice itself --- for instance via spatial translations, rotations or reflections --- or as time translation or time-reversal. The one-dimensional spacetime symmetries are the topic of \cref{app:frieze}. In what follows, we restrict attention to internal symmetries, which act directly on the quantum degrees of freedom rather than on the underlying spacetime. Concretely, we here assume the internal symmetries to be implemented \emph{globally} via a unitary \emph{on-site} (tensor product) representation of a group $G$. Crucially, the ground state projector commutes with the symmetry. The \emph{ground state subspace}, or \emph{ground space}, however, need not be trivial: even though the ground space may be one-dimensional, in which case it is \emph{symmetric}, it could also transform in a higher-dimensional representation of $G$, which signals \emph{spontaneous symmetry breaking} of $G$. In that case, a reference ground state is stabilised under the $G$-action by a maximal subgroup $H\subseteq G$, and one speaks of the \emph{symmetry-breaking pattern} $G\mapsto H$.

A gapped \emph{quantum phase of matter} can roughly be understood as an \emph{equivalence class} of gapped Hamiltonians whose properties vary smoothly with the Hamiltonian parameters. More precisely, two models are said to belong to the same phase when they can be interpolated by a smooth path of gapped, bounded-strength, local Hamiltonians that commute with the symmetry group $G$. From the point of view of quantum information theory, this is equivalent to the existence of a symmetric quantum circuit of sub-linear depth that maps their ground states into each other~\cite{Hastings2005b,Bravyi2006,Hastings2010}. Distinct gapped phases are separated by \emph{continuous phase transitions} characterised by a diverging correlation length and \emph{algebraic} decay of correlation functions whose universal scaling behaviour is captured by a \emph{conformal field theory}~\cite{DiFrancesco1997}. For a Hamiltonian with two controllable parameters, a fictitious phase diagram with three gapped phases is sketched below. The dotted line represents a Hamiltonian path interpolating between phases \Romannum{2} and \Romannum{3}. Along this path the gap closes at the phase boundary.
\begin{equation}
	\inputtikz{PhaseDiagram}\,\,.
\end{equation}
Built upon intuition about classical many-body systems, \emph{Landau's paradigm} has since long provided the unifying framework of phases and phase transition~\cite{Landau1937}. This framework states that phases of matter are in one-to-one correspondence with symmetry-breaking patterns of the symmetry. %\comment{The canonical classical example of this phenomenon is of course the transition between water and ice. The transition from the fluid to the solid phase is characterised by the breaking of the continuous translation symmetry of three-dimensional Euclidean space to the much smaller discrete symmetry of the ice lattice structure.}
Moreover, the paradigm posits that phase transitions can be diagnosed by local \emph{order parameters} that are charged under the symmetry group that acquire a non-vanishing vacuum expectation value in the symmetry-broken phase. Since Kadanoff and Ceva~\cite{Kadanoff1970} it has also been understood that \emph{symmetric} phases are analogously recognised via \emph{disorder parameters}.

The canonical example is the one-dimensional transverse field \emph{Ising model}, governed by the Hamiltonian~\cite{Degennes1963}
\begin{equation}\label{eq:TFIM}
	H = -J\sum_\msf i Z_\msf iZ_{\msf i+1} - \lambda \sum_\msf i X_\msf i,
\end{equation}
which possesses a global $\mbb Z_2$ \emph{spin-flip symmetry} generated by $\prod_\msf i X_\msf i$. In the \emph{ordered} or \emph{ferromagnetic} phase ($J > \lambda$), the spontaneous breaking of this symmetry is characterised by a non-vanishing expectation value of the local magnetisation $\la Z_\msf i \ra \neq 0$, which serves as the local order parameter.\footnote{One typically defines this order parameter by introducing a small longitudinal field $h\sum_\msf i Z_\msf i$ into the Hamiltonian, whose role is to select one of the two symmetry-broken ground states depending on the sign of $h$, which is then taken to zero, $h\mapsto 0$, after going to the thermodynamic limit $N\mapsto+\infty$.} In the \emph{disordered} or \emph{paramagnetic} phase ($\lambda > J$), the symmetry is restored and $\la Z_\msf i \ra = 0$, but now the disorder parameter $\la \prod_{\msf i=\msf k}^\msf l X_\msf i\ra$ attains a non-zero vacuum expectation value.

\bigskip\noindent Since its conception, quantum phases of matter that defy Landau's paradigm have emerged. Perhaps unsurprisingly, these \emph{new} phases are characterised by the distinct ways in which the symmetry can manifest itself in the entanglement patterns of the ground state~\cite{Bravyi2006,Schuch2011,Chen2010,Zeng2015,Wen2016,McGreevy2022,Chen2025}. Given the prominent role of some of these phases throughout this thesis, let us take some time here to introduce them and discuss some of their properties.

\bigskip\noindent In the presence of an on-site group symmetry, gapped \emph{symmetry-protected topological} or \emph{trivial} (SPT) phases are characterised by the uniqueness of the ground state on periodic boundary conditions and are adiabatically connected to product states, when the symmetry is explicitly broken~\cite{Chen2010,Chen2011}. On open boundary conditions they exhibit gapless modes that are exponentially localised near the boundaries, and which transform \emph{projectively} under the symmetry group. In one dimension, they can be completely described within the framework of one-dimensional tensor network states or \emph{matrix product states}, and we will revisit their classification and physical properties within this framework in \cref{sec:1Dphases}.

Unlike SPT phases, \emph{(intrinsic) topological phases} exist only in two or more spatial dimensions, and do not require the presence of a (physical) symmetry. They are characterised by a ground state degeneracy that depends solely on the topology of the spatial manifold and their gapped excitations exhibit \emph{anyonic} braiding statistics. A particular class of topological phases, which admit a microscopic lattice realisation in terms of \emph{string-net models}, form the topic of \cref{ch:topo}, and are also naturally described in terms of tensor networks. As shown below, these tensor networks host \emph{freely deformable} symmetry operators that act only on the \emph{entanglement degrees of freedom}. Such symmetries furnish representations of \emph{fusion categories}, thereby extending familiar notions from finite group and representation theory to algebraic structures whose multiplication law does not generally have an inverse.
\section{Outline of this thesis}
This thesis is divided into two parts: \hyperref[part:Foundations]{\emph{Foundations}} and \hyperref[part:Applications]{\emph{Applications}}. The latter entails the publications around which the thesis is centred and contains the research results obtained over the past five years. All of this work was carried out in collaboration with others. Each chapter begins with a brief summary of the main results, references the corresponding publication or preprint, and details my own contributions to the project. The research article is then reproduced in its published form or in its most recent version. One of these chapters is a set of lecture notes on tensor networks and as such does not contain new research results.

The first part on the other hand, consists of three chapters that serve both as an extended introduction to the topics underlying the research presented in the second part and as a guide to the relevant literature. Chapter \ref{ch:TN} introduces the framework of tensor networks, which constitute the primary technical tool employed throughout this thesis. Chapter \ref{ch:topo} provides an overview of two-dimensional non-chiral topological orders realised by string-net models and discusses their tensor network representations. Finally, chapter \ref{ch:GenSym} is devoted to the study of the corresponding one-lower-dimensional, viz.(1+1)-dimensional, \emph{edge theories} and their inherited \emph{generalised symmetry structures}, with an outlook towards the (2+1)-dimensional setting.

The thesis concludes with a summary of the main results, together with some general remarks and perspectives for future research.